\begin{abstract}
We revisit a negative result on joint-embedding predictive architectures (JEPA) for rotor dynamics and show that it was mostly structural: speed was observed directly, the latent prediction never reached the forecast, and constant commands made control authority unidentifiable. We rebuild the benchmark around the problem JEPA is meant for. A rotor governed by Euler turbomachinery work, with a lagged and delayed actuator and six hidden per-episode parameters, is observed only through a synthetic casing-microphone spectrogram with distractor tones, bursts and gain drift. A two-stage, action-conditioned JEPA world model learns from 3000 unlabelled episodes. Masked time--frequency latent prediction pre-trains the encoder, and a frozen-target action-conditioned predictor then learns the dynamics in latent space. It is compared under identical budgets with a reconstruction world model, supervised learning at 10--3000 labels, and a strong classical pipeline (harmonic tracker plus grey-box identification of the true ODE form). With a few labelled episodes, a frozen JEPA latent forecasts rotor speed several times better than supervised training from scratch. It degrades gracefully under nuisance levels at which the classical tracker fails, recovers control authority, and supports closed-loop set-point tracking by planning in latent space close to an oracle controller. The classical pipeline, which is given the exact equation, remains more accurate in distribution. The latent recovers state far better than it identifies the hidden machine parameters. All data are synthetic; nothing here validates an engine, sensor or controller.
\end{abstract}
